\chapter{Introduction}
\label{ch:intro}

\section{Bayesian model calibration}

A physical model predicts observations $y$ from a vector of parameters $\theta \in \R^d$.
\emph{Calibrating} the model means answering two questions: which parameters make it agree with
the data, and how well the data actually pin those parameters down. Bayesian model calibration
answers both at once, with the posterior distribution
\begin{equation}
  p(\theta \mid y) \;\propto\; \mathcal{L}(y \mid \theta)\,\pi(\theta),
  \label{eq:bayes}
\end{equation}
the likelihood of the data under the model, times a prior. In practice one works with its
logarithm,
\begin{equation}
  \log p(\theta) \;=\; \log \mathcal{L}(y \mid \theta) + \log \pi(\theta) + \text{const},
  \label{eq:logpost}
\end{equation}
which is only ever known up to that additive constant. The constant does not matter: the
posterior is unchanged by $\log p \to \log p + c$, and every algorithm in this manual is built so
that it never needs it.

The obstacle is cost. Each evaluation of $\log p$ runs the forward model --- a simulation that can
take minutes or hours --- and the algorithms that explore a posterior need thousands to millions of
evaluations. \textbf{slurmcmc} makes this tractable on a Slurm cluster in three ways: it evaluates
as many parameter points \emph{simultaneously} as an algorithm allows; it keeps a complete,
restartable record of every evaluation, so no expensive work is ever repeated; and, when even that
is not enough, it replaces most evaluations with a fast surrogate model while measuring whether the
surrogate can be trusted.

\section{The workflow}

A calibration typically goes through three stages, which is also how this manual is organised.

\begin{enumerate}[leftmargin=*, label=\textbf{\arabic*.}]
  \item \textbf{Optimization} (\cref{ch:optimization}). Minimize a loss that measures the misfit
        to data. This checks that the model \emph{can} represent the data at all, before any
        effort is spent on uncertainty, and it finds the region where the posterior lives, which
        is where MCMC should start.
  \item \textbf{MCMC} (\cref{ch:mcmc}). Sample the posterior with an ensemble Markov chain Monte
        Carlo algorithm whose walkers are evaluated in parallel on the cluster.
  \item \textbf{Surrogate-based MCMC} (\cref{ch:surrogate,ch:hybrid}). When a full MCMC is
        unaffordable, fit a fast surrogate to a limited number of expensive evaluations and sample
        the surrogate instead. The automated \emph{hybrid} pipeline iterates this until measured
        criteria say the surrogate posterior can be trusted.
\end{enumerate}

All three share the same machinery for dispatching evaluations to the cluster, recording them,
and restarting interrupted runs. \Cref{ch:parallel} describes it once, so the later chapters can
concentrate on the algorithms.

\section{What the package is built on}

slurmcmc stitches together established libraries rather than re-implementing them:

\begin{center}
\small
\begin{tabularx}{\linewidth}{@{}l L@{}}
  \toprule
  \textbf{library} & \textbf{role in slurmcmc} \\
  \midrule
  \href{https://github.com/facebookincubator/submitit}{\texttt{submitit}} &
    submits evaluations as Slurm job arrays (or local processes) and collects their results \\
  \href{https://github.com/facebookresearch/nevergrad}{\texttt{nevergrad}} &
    gradient-free population optimizers: Differential Evolution, Particle Swarm, and others \\
  \href{https://github.com/pytorch/botorch}{\texttt{botorch}} &
    Bayesian optimization with a Gaussian-process model \\
  \href{https://github.com/dfm/emcee}{\texttt{emcee}} &
    the affine-invariant ensemble MCMC sampler \\
  \href{https://scikit-learn.org}{\texttt{scikit-learn}} &
    the Gaussian-process and polynomial surrogate models of the hybrid pipeline \\
  \bottomrule
\end{tabularx}
\end{center}

and exposes four entry points:

\begin{center}
\small
\begin{tabularx}{\linewidth}{@{}l l L@{}}
  \toprule
  \textbf{entry point} & \textbf{module} & \textbf{purpose} \\
  \midrule
  \code{slurm_minimize}    & \code{slurmcmc.optimization} & parallel minimization of a loss (\cref{ch:optimization}) \\
  \code{slurm_mcmc}        & \code{slurmcmc.mcmc}         & ensemble MCMC with walkers evaluated on the cluster (\cref{ch:mcmc}) \\
  \code{slurm_mcmc_hybrid} & \code{slurmcmc.hybrid}       & automated surrogate-based MCMC (\cref{ch:hybrid}) \\
  \code{SlurmPool}         & \code{slurmcmc.slurm_utils}  & a cluster-backed \code{map}, used by all of the above and usable on its own (\cref{ch:parallel}) \\
  \bottomrule
\end{tabularx}
\end{center}

\section{Installation}

Install the package from the repository root; its core dependencies are pulled in automatically:
\begin{lstlisting}[style=plain]
pip install -e .
\end{lstlisting}
Bayesian optimization needs the \code{botorch} extra and the hybrid pipeline the \code{hybrid}
extra (scikit-learn). The \code{examples} and \code{test} extras add what the example scripts and
the test suite need. To install everything at once:
\begin{lstlisting}[style=plain]
pip install -e ".[botorch,hybrid,examples,test]"
\end{lstlisting}
Python~3.10 or newer is required. On a shared cluster without administrator rights, add
\code{--user} to install into your home directory.

\section{Running the tests}

\begin{lstlisting}[style=plain]
pytest -vv tests
\end{lstlisting}
Tests that need a real Slurm cluster are skipped automatically elsewhere; a mock-Slurm test file
exercises the Slurm code path without one.

\section{The examples behind this manual}

Every figure in this manual is produced by one of the example scripts in the \code{examples}
directory, so every result shown can be reproduced, and each script doubles as a starting point for
a real problem.

\begin{center}
\small
\begin{tabularx}{\linewidth}{@{}l L@{}}
  \toprule
  \textbf{script} & \textbf{shown in} \\
  \midrule
  \code{example_optimization.py}                   & \cref{sec:opt-example} \\
  \code{example_optimization_algorithms_comparison.py} & \cref{sec:opt-comparison} \\
  \code{example_mcmc.py}                           & \cref{sec:mcmc-example} \\
  \code{example_mcmc_and_mc_comparison.py}         & \cref{sec:mcmc-vs-mc} \\
  \code{example_mcmc_surrogate.py}                 & \cref{sec:surrogate-example} \\
  \code{example_mcmc_hybrid.py}                    & \cref{sec:hybrid-example} \\
  \bottomrule
\end{tabularx}
\end{center}
